\setcounter{section}{6}\refstepcounter{section}\addcontentsline{toc}{section}{7 Conclusion and Future Work}
\section*{7 Conclusion and Future Work}

\subsection*{7.1 Findings}

\begin{enumerate}
\item \textbf{(RQ1) Selection depends on the dataset.} Cost-weighted LR leads the
credit-g main matrix at F1 0.597; RF + SMOTE leads sick at 0.897. Section 6.2 supports
recall gains more strongly than F1 gains. No intervention improves every model on both
datasets.
\item \textbf{(RQ2) Accuracy needs minority-class evidence.} The reference rows in
Tables 6.1--6.2 illustrate the limitation of accuracy. Positive-class precision, recall
and F1, together with relevant error costs, reveal the decision trade-offs.
\item \textbf{(RQ3) One-hot encoding favours the tested LR pipeline; imputation
preserves sick records.} One-hot reaches F1 0.531 against 0.406 with ordinal codes on
credit-g LR; imputation retains all 3{,}772 sick records with the best AUC among the
three tested missing-data schemes (0.974). These conclusions are scoped to the tested
configurations; dataset differences and shared explanatory fields do not establish
causal mechanisms.
\item \textbf{(RQ4) The champions cut holdout cost and missed positives.} On the
holdout, credit-g cost falls to 0.570 from 0.887 for plain LR and 1.500 for
always-approve ZeroR, and the sick champion misses 3 of 69 positives with ten false
alarms. Section 6.4 bounds these results to internal evaluation because model
selection and final testing were not fully independent.
\end{enumerate}


\subsection*{7.2 Limitations}
These results describe internal comparisons on two datasets using one split seed. They do not establish performance on new populations. Retained duplicate sick records may appear in different validation folds. Configuration selection used cross validation over all records, so the overlapping holdout does not provide an independent test of that selection. Ordinary SMOTE can create artificial categorical representations, but we did not isolate their effect on performance. The ensemble extension and Apriori analysis are exploratory. We did not directly test the higher ensemble F1 against the main experiment winner. These limits leave uncertainty about generalisation and the reasons for model differences.

\subsection*{7.3 Future Work}
The next step is to test the selected configurations on independent data, or use nested cross validation with several seeds. A controlled comparison with SMOTENC could test whether preserving categorical values improves resampling. Threshold selection should follow the decision objective, including the official credit cost matrix where applicable. The Apriori rules also need checks on separate samples to see whether their support, confidence and lift remain stable.

